\section{Eigenvalues and normalized powers}
\label{sec:contraction_sums_power}

Throughout this section $H=\sum_{i\in\mathcal I}h_i$ on a finite-dimensional
complex Hilbert space, where $\mathcal I$ is finite and $0\le h_i\le1$, and
$N=|\mathcal I|+1$.

\begin{lemma}[The range of an eigenvalue]
    \label{thm:positive_sum_eigenvalue_range}
    \lean{Matrix.eigenvalue_nonneg_and_le_card_of_posSemidef_terms}
    \leanok
    If $H\Omega=E\Omega$ for a unit vector $\Omega$, then
    $0\le E\le|\mathcal I|$.
\end{lemma}

\begin{proof}
    \leanok
    Summing the bounds on the terms gives $0\le H\le|\mathcal I|\,1$, and
    $E=\langle\Omega,H\Omega\rangle$. This is the range of the ground energy
    used in~\cite[lines 283--284 and 607]{Landau2013PolynomialTime} to define
    the normalization $(1-E_0/N)^{-1}$.
\end{proof}

\begin{theorem}[Contraction of the excited vectors from a lower bound on the gap]
    \label{thm:physical_power_filter_lower_gap_cutdown}
    \lean{Matrix.norm_power_filter_apply_le_ratio_of_lower_gap}
    \leanok
    \uses{thm:spectral_power_ground_projection,thm:global_gap_ground_projection,
        thm:gaussian_gap_dichotomy}
    Let $\Omega$ be a unit vector, $E_0\in\mathbb R$ and $\Delta>0$ with
    $H\Omega=E_0\Omega$ and $H-E_0\,1\ge\Delta(1-P)$, where
    $P=|\Omega\rangle\langle\Omega|$. Put $T=(1-H/N)/(1-E_0/N)$ and
    \begin{align*}
        r=\frac{1-(E_0+\Delta)/N}{1-E_0/N}.
    \end{align*}
    For every unit vector $\psi$ orthogonal to $\Omega$ and every integer
    $m\ge0$,
    \begin{align*}
        \|T^m\psi\|                  & \le r^m, &
        |\langle\psi,T^m\psi\rangle| & \le r^m.
    \end{align*}
    This is equation e:cutdown of~\cite[line 603]{Landau2013PolynomialTime},
    with $\Delta$ any positive lower bound on the gap: $E_0+\Delta$ need not
    belong to the spectrum of $H$.
\end{theorem}

\begin{proof}
    \leanok
    The gap inequality at $\psi$ gives $\langle\psi,H\psi\rangle\ge E_0+\Delta$,
    and $H\le N-1$ gives $\langle\psi,H\psi\rangle\le N-1$. Hence
    $E_0+\Delta\le N-1$, and by the gap dichotomy the spectrum of $H$ lies in
    $\{E_0\}\cup[E_0+\Delta,N]$. Theorem~\ref{thm:spectral_power_ground_projection}
    at $E_1=E_0+\Delta$ and $M=N$, with
    Theorem~\ref{thm:global_gap_ground_projection} identifying its ground
    projection with $P$, gives $\|T^m-P\|\le r^m$, where
    $r=(N-E_0-\Delta)/(N-E_0)$. Since $P\psi=0$,
    $|\langle\psi,T^m\psi\rangle|\le\|T^m\psi\|=\|(T^m-P)\psi\|\le r^m$.
\end{proof}

\begin{lemma}[The affine filter normalized by the operator norm]
    \label{lem:full_ground_affine_filter_norm}
    \lean{Matrix.PosSemidef.normalized_affine_ground_filter_spec_of_norm}
    \leanok
    \uses{lem:full_ground_normalized_affine_filter}
    Let $H\ge0$ have eigenvalues $\mu<\nu$, let $\gamma>0$, and suppose that
    no point of the spectrum of $H$ other than $\mu$ lies below $\mu+\gamma$.
    Put $M=\|H\|$, $c=\gamma/(M-\mu)$, $K=(M-\mu)^{-1}(M-H)$, and let $G$ be
    the span of the eigenvectors of $H$ with eigenvalue at most $\mu$. Then
    $M>0$, $M-\mu>0$ and $0<c\le1$; moreover $Kx=x$ for $x\in G$, $K$ maps
    $G^\perp$ into itself, and $\|Kx\|^2\le(1-c)\|x\|^2$ for $x\in G^\perp$.
\end{lemma}

\begin{proof}
    \leanok
    The spectrum of $H$ lies in $[0,M]$, so $\mu\ge0$ and
    $\mu+\gamma\le\nu\le M$. Apply
    Lemma~\ref{lem:full_ground_normalized_affine_filter} with $E_1=\mu+\gamma$;
    its rate is $(E_1-\mu)/(M-\mu)=c$. This is the filter of the final step
    of~\cite[lines 1620--1643]{Arad2016RigorousRG} under its hypothesis (DG).
\end{proof}
